\section{总结与延伸}

本讲通过 Compound AI 的三步路径：1）捕获并压缩搜索轨迹，2）用 APAC 宪法治理多轮问答，3）通过 DSPy 学习复杂约束，在 1 小时内完整展示了从建模到工程落地的闭环。

这三部分之间不是孤立的：trace 模型用来指导 multi-turn agent 的问答流程，APAC 生成的结构化 Json 又是 DSPy 里 latent solver 的约束输入，最终的行动计划被送回 search trace 进行一致性检查，形成一个可审计的控制环。

\begin{table}[H]
\centering
\begin{tabular}{lp{5cm}p{6cm}}
\toprule
主题 & 关键机制 & 教学价值 \\
\midrule
搜索轨迹 & Trace-aware Transformer + DU Forer drop & 提高 data/parameter efficiency、增强可解释性 \\
多轮交互 & APAC Constitution + agent judge/DPO & 多轮提问具备 accuracy/proactivity/credibility \\
复杂约束 & Latent cost \(C(Y)\) + solver + green descent & 在非线性场景中嵌入可微反馈、兼顾制造约束 \\
\bottomrule
\end{tabular}
\caption{Compound AI 三步法的核心组件与价值}
\end{table}

\subsection{研究问题与扩展}

\paragraph{AI 教练层的版本管理}
当 search trace、APAC agent、DSPy solver 同时迭代时，容易出现 ``trace 参数与 agent 判断不同步''。建议把 trace ID、agent judge 模型版本、solver configuration 一起打包成 release 记录，在每次部署前跑自动化 test，确保新版 judge 仍能识别旧 trace，避免 unstable inference。

\paragraph{跨域泛化与 judge 校准}
当前的 agent judge 在 travel intent 任务上表现突出，但在医疗、金融等高风险领域可能对 goal 的理解截然不同。可行的扩展策略是训练 judge 的 mixture-of-experts 或引入 domain embedding 让 judge 自适应 new intent，同时把 judge reward 与 human satisfaction signal（NPS、回访率）对齐，而不是仅依赖 plan optimality。

\paragraph{多模态融合与辅助材料}
讲者提到 ``slides、notebooks、photos'' 是加强解释的利器。理想情况下，search trace 应该能够引用 slides 中的表格/公式（如 Soang 的 sequencing graph），当视频里出现关键 diagram 时自动截帧并插入 note，保持视觉与文本双重证据。

\begin{itemize}
    \item 如何进一步压缩 trace 的 token 数量以降低推理成本，同时保持多轮可解释性？
    \item agent judge 的泛化能力如何迁移到医疗、金融等高风险领域，而不仅仅是在 travel intent 里？
    \item 微分化的制造约束能否自动抽象成 latent representation，从而在新行业（例如芯片设计）间无缝迁移？
\end{itemize}

除了理论问题，这套 Compound AI 也要在 production 时持续沉淀：保持 trace/plan 的 audit log、把 DPO reward 与用户满意度对齐，并在不同 region 迭代 judge 模型，才能确保系统的长期健壮。

\subsection{拓展阅读}
\begin{itemize}
    \item ``Beyond A*: Better Planning with Transformers via Search Dynamics Bootstrapping''，用于理解搜索轨迹如何提升规划性能。
    \item ``Dualformer: Controllable Fast and Slow Thinking by Learning with Randomized Reasoning Traces''，对应 trace drop 与可控推理速度。
    \item ``Composing Global Optimizers to Reasoning Tasks via Algebraic Objects in Neural Nets''，讨论神经网络中的组合式全局优化。
    \item ``SurCo: Learning Linear Surrogates For Combinatorial Nonlinear Optimization Problems''，展示组合优化问题中的 surrogate 学习路径。
    \item APAC Constitution（2024），内部白皮书，详细定义了 Accuracy/Proactivity/Adaptivity/Credibility 的评价方法。
\end{itemize}

\subsection{本章小结}
Compound AI 的三步法不是孤立的实验，而是面向 production 的组合：把 visual slides、time-stamped frames、trace logs 融合进同一条 pipeline 里，让观众不仅理解方法原理，还能追踪到具体的工程实现。

\end{document}
